\cotefile{90}
\scanpage{2}
\deleted{\[f(gh)=\sum f_i(g)\varphi_i(h),\]}
\margin{\uline{$M$}}
\deleted{\illegible{1}}
\deleted{\unsure{fini et projectif|plat et projectif}}

\uline{Lemme 1} Soit $A$ un anneau de Dedekind, $M$ un $A$-module plat. \deleted{Conditions équivalentes} Pour $\forall$ sous-module $N$ de $M$, soit $\tilde N$ l'image inverse dans $M$ du sous-module de Torsion de $M/N$. Donc $\tilde N/N$ est un module de torsion.

\unsure{Cond.} équivalentes a) b) c) d)

a) $\forall\,N$ de type fini, $\tilde N/N$ est de longueur finie, i.e. $\tilde N$ est de t.f.

\margin{$\unsure{\Uparrow|\Updownarrow}$}
\deleted{\unsure{$A$ est un corps,}}
b) \inserted{Si $A$ n'est pas un corps,} Il n'existe pas de sous-module de $M$ isomorphe à :
\[
\inserted{A^{(p)}\unsure{=}\bigcup_n p^{-n}=\Gamma(S-\{p\},\underline O_S)\quad(S=\operatorname{Spec}A),}
\deleted{\unsure{K_p}}\quad p\text{ est un idéal maximal}
\]
de $A$ \deleted{$A$ est un corps.}

\margin{$\Updownarrow$}
\deleted{\unsure{Pour tout}}
c) \inserted{Si $A$ n'est pas un corps,} Pour $\forall$ idéal maximal \deleted{\unsure{$p$}} de $A$, $M_p$ ne contient pas de \unsure{$A$-module} Isom. à $K$.

d) $M=\varinjlim M_i$, les $M_i\subset M$ de type fini et ``\unsure{saturés}'', i.e. \inserted{$M_i\longrightarrow M_j$ \unsure{ts que|tels que} $M_i$ fact direct dans $M_j$.}

Pour que ces conditions soient \deleted{\unsure{vraies}}, il suffit que l'on ait :

e) Les formes linéaires sur $M$ séparent, i.e.
\[
M\hookrightarrow A^I\qquad(I\text{ un ensemble})
\]
[condition vérifiée si $M$ contenu ds un libre, \deleted{\unsure{par}} i.e. $M$ projectif]

C'est aussi nécessaire si $M$ a un ens. dénombrable de générateurs.
\mspage{1}
